\documentclass[11pt,a4paper]{article}

\usepackage[margin=1in]{geometry}
\usepackage{booktabs}
\usepackage{siunitx}
\usepackage{amsmath}
\usepackage{graphicx}
\usepackage{hyperref}
\usepackage{longtable}
\usepackage{xcolor}

\sisetup{detect-all=true}
\graphicspath{{../figures/Detector_performance_analysis_Rn222_repetition_comparison/}{figures/Detector_performance_analysis_Rn222_repetition_comparison/}}

\title{Photon Simulation Model and Rn-222 Detector-Response Results}
\author{Theranostics simulation workspace}
\date{\today}

\begin{document}


\usepackage[margin=1in]{geometry}
\usepackage{booktabs}
\usepackage{siunitx}
\usepackage{amsmath}
\usepackage{graphicx}
\usepackage{hyperref}
\usepackage{longtable}
\usepackage{xcolor}

\sisetup{detect-all=true}
\graphicspath{{../figures/Detector_performance_analysis_Rn222_repetition_comparison/}{figures/Detector_performance_analysis_Rn222_repetition_comparison/}}

\title{Photon Simulation Model and Rn-222 Detector-Response Results}
\author{Theranostics simulation workspace}
\date{\today}

\begin{document}
\maketitle

\begin{abstract}
This note summarizes a photon-transport model for Rn-222 decay-chain
activity near a thin glue layer coupled to a water phantom.  The current
results show that changing the Rn-222 spatial distribution has a modest
effect on total photon-detector entries, but a strong effect on the
relative loading of the six detector panels.  Moving activity from the
glue layer into water reduces the dominant positive-\(y\) response and
increases the negative-\(y\), \(x\)-panel, and \(z\)-panel responses.
\end{abstract}

\section{Simulation Objective}

The simulation models photon signatures from Rn-222 decay-chain activity
near a thin glue layer coupled to a water phantom.  The practical question
is whether photon-detector hit patterns can distinguish activity retained
near the glue from activity that has diffused or been placed elsewhere in
the water volume.  The analysis considers detector-entry counts, panel
profiles, energy and time spectra, source attribution, and simple
back-projection diagnostics.

\section{Geometry and Materials}

The model contains:

\begin{itemize}
  \item a \(2 \times 2 \times 2~\si{m}\) air world;
  \item a water box with half-lengths \(x=\SI{10}{cm}\), \(y=\SI{10}{cm}\),
        and \(z=\SI{75}{cm}\);
  \item a cyanoacrylate glue slab with half-lengths \(x=\SI{2}{cm}\),
        \(y=\SI{0.5}{mm}\), and \(z=\SI{2}{cm}\), centered at
        \(y=\SI{100.5}{mm}\);
  \item six silicon photon-detector plates surrounding the water box at
        \(\SI{20}{cm}\) from the water surface, with an additional
        \(\SI{5}{cm}\) lateral padding.
\end{itemize}

The glue material is represented as cyanoacrylate with density
\(\SI{1.07}{g.cm^{-3}}\) and elemental composition \(C_6H_7NO_2\).  The
water and silicon detector materials use standard reference material
definitions.

\section{Physics and Source Model}

The transport model includes low-energy electromagnetic interactions,
radioactive decay-chain evolution, and step limiting.  The production-cut
energy range is set to \(\SI{100}{eV}\) to \(\SI{1}{GeV}\), with a default
cut value of \(\SI{1}{nm}\).  Primaries are generated as stationary
Rn-222 ions:
\[
  Z=86, \qquad A=222, \qquad E=\SI{0}{MeV}.
\]

Rn diffusion is represented by sampling a random direction and a
decay-time-dependent displacement scale
\[
  \Delta r \sim \left| \mathcal{N}\left(0, \sqrt{6Dt}\right) \right|,
\]
with \(D=\SI{1.12e-3}{mm^2.s^{-1}}\) for radon.  The glue slab side walls
reflect the sampled displacement in \(x\) and \(z\).  Motion out of the top
of the glue is reflected, while motion below the bottom glue boundary is
logged as an Rn-222 glue-exit event.  The exit log records event ID, track
ID, estimated exit time, decay time, start position, exit position, and
final sampled position.

\section{Repeated Rn-222 Campaign}

The repeated-run campaign contains six source configurations, each with ten
independent repetitions and \(10{,}000\) simulated primary Rn-222 events per
repetition.  Each repetition uses a unique random seed.

\begin{table}[h]
\centering
\caption{Rn-222 source configurations in the repeated campaign.  The glue
source is a \(20 \times 1 \times 20~\si{mm}\) half-size volume centered at
\((0,101,0)~\si{mm}\).  The diffuse water source is the full water box.
The localized water source is a \(2.5~\si{cm}\) half-size cube centered at
\((0,0,50)~\si{cm}\).}
\begin{tabular}{lll}
\toprule
Case & Glue fraction & Additional Rn-222 source \\
\midrule
baseline & 1.0 & none \\
diff10 & 0.9 & 0.1 diffuse in water box \\
diff50 & 0.5 & 0.5 diffuse in water box \\
diff90 & 0.1 & 0.9 diffuse in water box \\
loc40 & 0.6 & 0.4 localized in water cube \\
loc60 & 0.4 & 0.6 localized in water cube \\
\bottomrule
\end{tabular}
\end{table}

\section{Analysis Definitions}

The primary analysis records gamma entries and exits at the water volume,
glue volume, and detector plates.  The detector-response metrics in this
note use gamma tracks entering a photon detector plate:
\[
  \mathrm{particle}=\gamma, \qquad
  \mathrm{boundary}=\mathrm{enter}, \qquad
  \mathrm{volume}=\mathrm{photon~detector}.
\]
Counts are therefore detector-entry counts, not all emitted photons.  The
analysis also forms unique detected gamma-track counts by de-duplicating
within each event.  Profile and heatmap comparisons use baseline-subtracted
detector-entry histograms, with repeated runs providing the run-to-run
uncertainty.

\section{Results}

\subsection{Total Detector Entries}

Table~\ref{tab:total-response} summarizes the total detector entries per
\(10{,}000\) primary Rn-222 events.  The total response decreases from the
baseline to high water-distributed cases, but the change is small compared
with the much larger redistribution among detector panels shown below.

\begin{table}[h]
\centering
\caption{Mean detector-entry counts across ten repetitions.  Uncertainties
are sample standard deviations across repetitions.}
\label{tab:total-response}
\begin{tabular}{lccc}
\toprule
Case & Detector entries & Unique gamma tracks & Entries per primary event \\
\midrule
baseline & \(12500 \pm 82\) & \(12415 \pm 82\) & \(1.250 \pm 0.008\) \\
diff10 & \(12431 \pm 100\) & \(12349 \pm 97\) & \(1.243 \pm 0.010\) \\
diff50 & \(12300 \pm 130\) & \(12218 \pm 128\) & \(1.230 \pm 0.013\) \\
diff90 & \(12057 \pm 86\) & \(11981 \pm 87\) & \(1.206 \pm 0.009\) \\
loc40 & \(12174 \pm 79\) & \(12096 \pm 80\) & \(1.217 \pm 0.008\) \\
loc60 & \(12134 \pm 81\) & \(12058 \pm 80\) & \(1.213 \pm 0.008\) \\
\bottomrule
\end{tabular}
\end{table}

\subsection{Detector-Panel Redistribution}

Table~\ref{tab:panel-response} gives mean detector-entry counts by panel.
The baseline case is dominated by the positive-\(y\) plate because the glue
source lies immediately above the water box.  As more Rn-222 is assigned to
water, the positive-\(y\) excess is reduced and the negative-\(y\), \(x\),
and \(z\) panels receive larger shares.  The localized water cases retain
the same qualitative behavior but produce asymmetric \(z\)-panel responses
because the localized cube is placed at positive \(z\).

\begin{table}[h]
\centering
\caption{Mean detector-entry counts by detector panel across ten repetitions.}
\label{tab:panel-response}
\begin{tabular}{lrrrrrr}
\toprule
Case & NegX & PosX & NegY & PosY & NegZ & PosZ \\
\midrule
baseline & 2570 & 2594 & 1420 & 5882 & 18 & 16 \\
diff10 & 2622 & 2632 & 1550 & 5552 & 35 & 40 \\
diff50 & 2740 & 2743 & 2143 & 4412 & 133 & 130 \\
diff90 & 2851 & 2851 & 2736 & 3178 & 223 & 218 \\
loc40 & 2694 & 2683 & 1986 & 4663 & 11 & 136 \\
loc60 & 2741 & 2748 & 2315 & 4117 & 9 & 204 \\
\bottomrule
\end{tabular}
\end{table}

\begin{figure}[h]
\centering
\includegraphics[width=\textwidth]{y_panel_baseline_subtracted_z_profiles.png}
\caption{Baseline-subtracted \(z\)-profiles for the positive-\(y\) and
negative-\(y\) detector panels.  Moving Rn-222 from the glue layer into
water changes the panel balance and produces a measurable residual
structure along the detector length.}
\label{fig:y-profile}
\end{figure}

\begin{figure}[h]
\centering
\includegraphics[width=\textwidth]{y_panel_residual_significance_profiles.png}
\caption{Residual-significance profiles for the \(y\)-panel comparison.
The repeated runs are used to estimate the uncertainty on the
candidate-minus-baseline profile.}
\label{fig:y-significance}
\end{figure}

\begin{figure}[h]
\centering
\includegraphics[width=\textwidth]{y_panel_2d_residual_maps_loc60.png}
\caption{Two-dimensional \(y\)-panel residual maps for the localized
water-source case with 60\% of the activity in the localized water region.}
\label{fig:y-map}
\end{figure}

\begin{figure}[h]
\centering
\includegraphics[width=\textwidth]{x_panel_yz_residual_maps_loc40.png}
\caption{\(x\)-panel \(y\)-\(z\) residual maps for the localized water-source
case with 40\% of the activity in the localized water region.}
\label{fig:x-map}
\end{figure}

\begin{figure}[h]
\centering
\includegraphics[width=\textwidth]{y_panel_asymmetry_residual.png}
\caption{Positive-\(y\) minus negative-\(y\) asymmetry residual.  This
diagnostic emphasizes the directional shift in detector response as
activity leaves the glue-dominated source configuration.}
\label{fig:y-asymmetry}
\end{figure}

\subsection{Isotope Attribution}

Across all Rn-222 configurations, detected photons are dominated by
late-chain gamma emitters.  In the baseline case, approximately
\(\SI{60.5}{\percent}\) of detector entries are attributed to Po-214,
\(\SI{37.1}{\percent}\) to Bi-214, and \(\SI{2.2}{\percent}\) to Bi-210.
Pb-214 and other isotopes are present at the sub-percent level in the
detector-entry stream.  The isotope mix remains nearly unchanged across
the source-placement cases, so the strongest discriminant in the current
analysis is spatial detector response rather than a change in gamma-line
parentage.

\subsection{Timing}

Cumulative time cuts show how many detector entries would be available as
a function of observation time.  For the baseline case, the mean cumulative detector
entries are \(29.5\) at 1 hour, \(106.3\) at 2 hours, \(977.8\) at 12 hours,
\(1933.3\) at 1 day, \(3639.8\) at 2 days, \(7217.8\) at 5 days, and
\(10204.6\) at 10 days, per \(10{,}000\) simulated primaries.  The other
source configurations follow similar time evolution with lower totals in
the high-water-fraction cases.

\begin{figure}[h]
\centering
\includegraphics[width=\textwidth]{time_window_integrated_residual_vs_time.png}
\caption{Integrated baseline-subtracted residuals as a function of
cumulative observation window.}
\label{fig:time-integrated}
\end{figure}

\begin{figure}[h]
\centering
\includegraphics[width=\textwidth]{time_window_max_zscore_vs_time.png}
\caption{Maximum residual significance as a function of cumulative
observation window.}
\label{fig:time-zscore}
\end{figure}

\subsection{Rn-222 Glue Exit}

The dedicated Rn-222 exit log records how often sampled Rn-222 motion
crosses the bottom of the glue slab before decay.  In the all-glue baseline
case, \(4593 \pm 50\) Rn-222 tracks per \(10{,}000\) primaries exit the glue
in the logged downward direction.  The mean repetition-level median exit
time is \(4.81 \pm 0.13\) hours, while the mean 90th-percentile exit time is
\(32.64 \pm 1.27\) hours.  As expected, the number of glue-exit records
scales down with the simulated fraction of Rn-222 initially assigned to the
glue layer: \(4094 \pm 41\) for diff10, \(2278 \pm 46\) for diff50, and
\(456 \pm 18\) for diff90.

\section{Interpretation}

The current simulation and analysis support three conclusions.
First, the detector system is sensitive to source placement through panel
balance: total counts alone are a weak discriminator, but the relative
response of \(+y\), \(-y\), \(x\), and \(z\) plates changes systematically.
Second, the dominant detected gamma population comes from Po-214 and
Bi-214, so isotope attribution mainly confirms the expected Rn-222 decay
chain rather than separating the spatial hypotheses.  Third, the logged
diffusion model predicts substantial Rn-222 movement from the glue into
water on hour-scale times, which is consistent with the need for cumulative
time-window analysis and activity-normalized projections.

\section{Limitations and Next Steps}

The current tabulated analysis treats detector entries as ideal photon hits.  It
does not yet include detector efficiency, energy resolution, coincidence
logic, electronic thresholds, or measured background.  The back-projection
analysis estimates source-track proxies from detector-entry position and
momentum, but a single gamma constrains a line rather than a unique source
point.  A stronger reconstruction study should combine detector response
maps with a likelihood model over source distributions and should include
the detector-response model needed for comparison with data.

\end{document}

\end{document}
